\documentclass[11pt]{article}
\usepackage[T1]{fontenc}
\usepackage[utf8]{inputenc}
\usepackage{amsmath,amssymb,amsthm}
\usepackage[margin=0.9in]{geometry}
\usepackage{hyperref}
\setlength{\emergencystretch}{2em}
\hypersetup{hidelinks}

\newtheorem{theorem}{Theorem}
\newtheorem{lemma}[theorem]{Lemma}

\newcommand{\KG}{\mathrm{KG}}
\newcommand{\KGR}{R^{\mathrm{KG}}_{3}(3,3)}

\title{A Two-Coloring of the Edges of $\KG(11,3)$\\
with No Monochromatic Triangle}
\author{Xiaojian Zeng\\
PBC School of Finance, Tsinghua University, Beijing 100190, China\\
\texttt{xiaojianzeng@foxmail.com}}
\date{}

\begin{document}
\maketitle

\begin{abstract}
We give an explicit red--blue coloring of all $4{,}620$ edges of the
three-uniform Kneser graph $\KG(11,3)$ and exhaustively verify that none of
its $15{,}400$ triangles is monochromatic.  It follows that
$R^{\mathrm{KG}}_{3}(3,3)\geq 12$, improving the previously published lower
bound $11$.  Combined with the known upper bound $13$, this leaves the exact
value in $\{12,13\}$.  In particular, this local result does not decide the
$\KG(12,3)$ case, so it does not resolve the original exact-value problem.
The coloring and two verifiers using different enumerations are supplied as
machine-readable ancillary files.
\end{abstract}

\section{Statement and context}

For $n\geq 6$, the Kneser graph $\KG(n,3)$ has the three-element subsets of
$[n]=\{1,\ldots,n\}$ as vertices, with two vertices adjacent when they are
disjoint.  Following Heath, McCourt, Parker, Schwieder, and
Zerbib~\cite{HeathEtAl2025}, let $\KGR$ be the least $n$ such that every red--blue
edge-coloring of $\KG(n,3)$ contains a monochromatic triangle.  They proved
\[
  11\leq \KGR\leq 13
\]
and encoded the avoidance of monochromatic triangles as a satisfiability
problem~\cite[Theorems 11--12 and Section 3.3]{HeathEtAl2025}.

An edge-coloring without a monochromatic triangle is equivalent to a cover
of the edge set by two triangle-free spanning subgraphs: the color classes
give such a cover, while an edge in an overlapping cover may be assigned to
either covering graph.  Gujgiczer, Marits, and Ozeki used this formulation
and explicitly recorded the $\KG(11,3)$ and $\KG(12,3)$ triangle-free-cover
cases as unresolved~\cite[Section 4]{GujgiczerEtAl2026}.  Their separate
question about covers by two graphs of fractional chromatic number at most
$2.5$ is stronger and is not used here.

The contribution of this note is restricted to the following finite case.

\begin{theorem}\label{thm:main}
The graph $\KG(11,3)$ has a red--blue edge-coloring with no monochromatic
triangle.  Consequently,
\[
  R^{\mathrm{KG}}_{3}(3,3)\geq 12.
\]
\end{theorem}

Relative to the two sources above, Theorem~\ref{thm:main} decides the
$n=11$ case on the colorable side.  The literature comparison in this note
is bounded to the cited records and is not a claim of priority.

\section{Certificate proof}

We first record the finite objects that the certificate must cover.  There
are
\[
  \binom{11}{3}=165
\]
vertices.  Every vertex is disjoint from $\binom{8}{3}=56$ other vertices,
so the number of unordered edges is
\begin{equation}\label{eq:edges}
  \frac12\binom{11}{3}\binom{8}{3}=4{,}620.
\end{equation}
Three vertices form a triangle precisely when they are three pairwise
disjoint triples.  Equivalently, they are an unordered partition of a
nine-element subset of $[11]$ into three triples.  Thus the number of
triangles is
\begin{equation}\label{eq:triangles}
  \binom{11}{9}\frac{9!}{{(3!)}^3 3!}
  =55\cdot 280=15{,}400.
\end{equation}

The ancillary text file \verb|n11_coloring.txt| has one data line for every
edge.  A line
\[
  a_1,a_2,a_3\quad b_1,b_2,b_3\quad C
\]
assigns $C\in\{\mathrm{R},\mathrm{B}\}$ to the edge joining the displayed
sorted triples.  The primary verifier performs the following exact checks.

\begin{enumerate}
  \item It parses every endpoint, checks that it is a three-subset of
  $[11]$, checks disjointness, and rejects duplicate edges or invalid
  colors.
  \item Independently of the lines present in the certificate, it constructs
  all $165$ vertices and all unordered pairs of disjoint vertices.  It
  requires set equality between these expected edges and the parsed edges;
  the common cardinality is $4{,}620$ as in~\eqref{eq:edges}.
  \item It enumerates every unordered triple of graph vertices.  A triple is
  tested exactly when the union of its three three-subsets has size nine,
  which is equivalent to pairwise disjointness.  For each such triangle, it
  looks up its three edge colors and rejects equality of all three.  The
  number tested is required to be $15{,}400$ as in
  Equation~\eqref{eq:triangles}.
\end{enumerate}

The verifier accepts the supplied file.  It finds $2{,}308$ red edges and
$2{,}312$ blue edges.  Among all triangles, $7{,}720$ have one red and two
blue edges, while $7{,}680$ have two red and one blue edge.  Hence there are
no all-red or all-blue triangles.

\begin{proof}[Proof of Theorem~\ref{thm:main}]
The coverage check makes the parsed certificate a total map from
$E(\KG(11,3))$ to $\{\mathrm{R},\mathrm{B}\}$.  By
Equation~\eqref{eq:triangles}, the exhaustive triangle check considers every
triangle of the graph, and the reported two nonmonochromatic patterns
account for all $15{,}400$ of them.  The certificate therefore supplies the
claimed coloring.

The defining property is monotone in $n$, because $\KG(m,3)$ is the
subgraph of $\KG(n,3)$ induced by triples contained in $[m]$ whenever
$m\leq n$.  A coloring without a monochromatic triangle at $n=11$ therefore
implies that the least forcing value is greater than $11$, proving
$\KGR\geq 12$.
\end{proof}

\section{Independent replay and reproducibility}

A second verifier was written without importing the primary checker.  It
generates each edge by choosing its six-point union and then an unordered
partition of that union into two triples.  It generates each triangle by
choosing its nine-point union and then an unordered partition into three
triples.  These routes give, independently,
\[
  \binom{11}{6}\frac{\binom63}{2}=4{,}620,
  \qquad
  \binom{11}{9}\frac{9!}{{(3!)}^3 3!}=15{,}400.
\]
It obtains the same edge-color and triangle-pattern counts stated above.
As negative controls, it rejects a certificate with one edge deleted and
detects a monochromatic triangle after the three edges of a generated
triangle are forced red.

The two verification routes use integer set operations only; neither uses
floating-point arithmetic, randomized sampling, or a SAT solver.  A SAT
solver was used to discover the assignment, but its output is outside the
trust boundary: the explicit coloring is checked from scratch.  This
separation of discovery from exact certificate checking is also used in
other finite combinatorial work; for a directly verified example involving
stored graph colorings and independent reconstruction, see
Zeng et~al.~\cite[Sections 3 and 6]{202608.1658}.

The ancillary material contains the coloring, the primary standard-library
checker, the independently written audit verifier, and its recorded output.
Running
\begin{center}
  \verb|python3 check_n11_coloring.py|
\end{center}
from the directory containing the coloring and primary checker reproduces
the decisive coverage and triangle checks.  The independent verifier is
run from the project root as
\begin{center}
  \verb|python3 audit/referee_verify.py|.
\end{center}

\section{Scope and limitations}

Combining Theorem~\ref{thm:main} with the upper bound of Heath et~al.\ gives
only
\[
  12\leq R^{\mathrm{KG}}_{3}(3,3)\leq 13.
\]
Nothing in the $n=11$ certificate determines whether every two-coloring of
$\KG(12,3)$ has a monochromatic triangle.  An exact decision at $n=12$
would require either a similarly complete coloring or a fully checkable
refutation of the corresponding finite constraints.  Thus the original
exact-value problem remains unresolved by this note.

\paragraph{Computational and AI assistance.}
OpenAI Codex assisted with the literature search, SAT formulation and
witness discovery, checker development, adversarial audit organization, and
initial manuscript preparation.  A SAT solver supplied the candidate
coloring.  Neither solver status nor language-model output is used as
mathematical evidence: the result rests on the released finite coloring and
the exhaustive exact verifiers described above.  No proof assistant was
used.

\bibliographystyle{alphaurl}
\bibliography{references}

\end{document}
